\section{把战果写进簿册：金的行政选择}

\subsection{接收旧官署，而不是从零画出州县}

金进入华北时，面对的是辽、宋都曾经营过的城市和制度遗产。道路已有名称，仓库已有看守者，田土和户口在不同层级已有登记，寺院和市场也形成自己的关系网。新政权可以撤换显眼的官员，却不能在一天内替换所有懂得地方账目、粮仓钥匙和诉讼习惯的人。因此，沿用、改名、重组和再征发常常同时发生。把这种过程简单说成“汉化”，会抹平金廷的主动选择，也会假定旧制度能够无摩擦地继续运转。

《金史》的百官、地理与食货诸志保留了制度名称和后来的定制，适合追问中枢希望建立什么；它们不是每一州县的工作日记。诏令写出应当怎样处理，反复申令反而可能表示地方不能顺利照办。墓志所载的迁徙、任官和家族捐施能给出一些具体地点，却多出于精英纪念的目的。都城遗址、边城遗存、钱币与窑址能改变我们对空间和生产的认识，也不能替沉默的大多数发言。阅读金的行政，正要让这些材料相互约束，而不是把任何一种材料充当全国的缩影。

\subsection{猛安谋克与州县：两套语言进入同一地方}

猛安谋克常被当作解释金朝的钥匙。它确实将军事编制、部众关系和供给责任连接起来，使金廷能在扩张时维持一支具有自身组织传统的力量。但它不是一只覆盖全境、职能恒定的抽屉。金统治的地域从东北延伸到华北和河南，人口、生产与既有行政传统相差很大；同一个官名或编制进入不同地方，可能承担不同程度的军役、土地、户籍和治安功能。

州县行政同样不能被想成不变的宋制残余。金廷必须借助熟悉地方文书的人办理赋税、司法与运输，又要避免这些网络成为独立的资源中心。官员的调动、考课和监督，是为了把中都的判断送到各路；实际执行仍取决于道路通畅、库存、季节和地方能否承担。所谓中央集权并不是一根从宫门伸到村社的直线，而是一系列会被渡口、荒年、驻军和人情关系折弯的传递。

一一四九年完颜亮即位，使中枢权力重新集中。此后迁都中都，更清楚地把政治中心放在华北交通网内。迁都并非只为接近南宋。中都附近可以联结燕云、河北、山东、河南与东北的线路，朝廷由此更容易获得各地消息，也更容易把都城消费、工程和驻军的负担推向周围地区。新中心带来的不是统一速度，而是更密集的协调需求。

\subsection{中都的城门之后}

都城考古所见的城垣、街道、作坊与宫城遗迹，使中都不再只是史书里一次迁都的名词。遗址告诉我们哪些空间曾被营建、修补或使用，却不自动复原每条街的居民和每个工匠的生活。营建一座都城需要木料、砖瓦、金属、运输和劳动力；这些物资如何取得、由谁承担、付给谁多少，往往比工程完成的日期更难确证。正史的宏大叙述容易将工程写成君主意志，物质遗存则提醒我们有无数被征调、雇佣或依附于作坊的人未被逐名记录。

都城还是分类的机器。官员、军人、商人、僧道、工匠、使节和被迁徙的家户进入城门时，可能被不同衙门以不同方式登记。分类有助于征税、治安和供给，却不等于人们只具有一种身份。一名军户也可能经营小规模交易，一座寺院既是礼拜地也可能保存土地和捐施的文书，家族的姻亲关系会跨越官署所画的界限。我们不应把这些重叠浪漫化为自由流动；它们也会在征役、诉讼和粮价上升时成为风险。

\subsection{完颜亮的南侵：命令走得比粮食快}

一一五六年金廷为南侵集结物资。制度文本和本纪可以让我们看见动员的方向，不能把它们换算为一张精确的仓储清单。大型远征要经过的道路、河流、船只、马匹、饲料、军器和民夫相互牵连；任何一项迟滞，前线的计划都会改变。南宋与金的战报都倾向于将对手的困难或本方的功劳组织成胜败故事，故本书不采用未经交叉确认的兵数和伤亡。

一一六一年渡江受挫、完颜亮被杀，完颜雍即位。这不是一场战役足以解释的政权更替。都城政治、远征消耗、军中不满与南方水路的限制共同压缩了选择。长江并非地图上一条等待跨越的线：渡河需要集中船只，集中船只要有岸上的仓储与护卫，而江南一侧的防御又能利用水网和地方供给。战役结果重新安排了中枢的资源，但它没有让边界上的征发、贸易和人口流动停止。

\subsection{世宗时期的修复不是静止}

世宗的统治常被概括为“守成”或“盛世”。这样的总评容易遮住修复实际意味着什么：停战后要整顿军队、重建仓储、处理地方官与军政编制、调节税粮和都城供给。政策能够改善某些环节，却不能同时解决所有地区的水旱、流民、债务和边防。金朝能在南北之间维持较长时期的和平，部分由于双方都承受继续大规模战争的成本；这是一种需要不断支付的均衡，而非天然稳定。

行政的得失也在这里显露。直接治理使中枢能够更准确地调用华北的人力与财物，并将新的官僚和城市网络连接起来；成本则是地方对军费、差役和工程的承受会被不断测试。短期内，统一的命令有助于恢复秩序；中期看，反复的征调可能消耗正要复耕和重建的家庭。金后来面对蒙古攻势时，必须依赖的正是这些已被长期动员的道路、粮区和地方网络。统治所解决的问题，也留下了下一次危机的压力。
